\chapter{Memory Map and Bus Access}
\label{ch:memory-map}

\section{The Device Segment}

The video device occupies one segment of the CPU's address space, 65,536 words long. The segment is chosen by the
system; the reference virtual machine maps the device at segment \texttt{0x000C}, so a register at offset
\texttt{0x0002} is at the CPU address \texttt{0x000C\_0002}. This manual gives every address as an offset from the base
of the segment.

\begin{table}[H]
	\centering
	\begin{tabular}{llp{6.5cm}}
		\toprule
		\textbf{Offsets}                 & \textbf{Region} & \textbf{Notes}                                                                                           \\
		\midrule
		\texttt{0x0000}--\texttt{0x00FF} & Registers       & Offsets \texttt{0x0000}--\texttt{0x0018} are defined (Chapter~\ref{ch:registers}); the rest are open bus \\
		\texttt{0x0100}--\texttt{0x5FFF} & Unmapped        & Open bus                                                                                                 \\
		\texttt{0x6000}--\texttt{0x60FF} & Palette         & 256 entries, one word each (Section~\ref{sec:palette})                                                   \\
		\texttt{0x6100}--\texttt{0x7FFF} & Unmapped        & Open bus; see Section~\ref{sec:palette-edge}                                                             \\
		\texttt{0x8000}--\texttt{0xFFFF} & Video memory    & 32,768 words of tile data and tilemaps                                                                   \\
		\bottomrule
	\end{tabular}
	\caption{Video Device Memory Map}
	\label{tab:memory-map}
\end{table}

\section{Registers, Palette and Video Memory}

The \textbf{registers} control the device. They are described in Chapter~\ref{ch:registers}.

The \textbf{palette} holds 256 colours in a small, fast memory that the renderer reads once per pixel. It can be read
and written by the CPU at any time.

The \textbf{video memory} (VRAM) holds everything else: tile data, tilemaps and, when they are specified, sprite tile
data. It is 32,768 words, which is 64 KB. The device does not interpret VRAM in any fixed layout. Each background has a
base address register for its tilemap and another for its tile data, and any VRAM range may hold either. The only
constraint is that ranges the renderer reads must not overlap ranges the program expects to stay untouched.

An address in VRAM is written as its offset in the segment, so VRAM begins at \texttt{0x8000}. The base address
registers hold addresses in this form: a tilemap at the start of VRAM has the base address \texttt{0x8000}, not
\texttt{0x0000}.

\section{Access Rules}

A read of a register, a palette entry or a VRAM word returns the stored value. A write stores the value. A read of an
open bus offset returns zero and a write to one is ignored. The device never inserts a wait state: every access
completes in the cycle it is made, whether or not it takes effect (Section~\ref{sec:memory-vram-access}).

The Status register (offset \texttt{0x0017}) and the Vertical Line register (offset \texttt{0x0018}) are read-only; a
write to either is ignored.

\subsection{Access During Rendering}
\label{sec:memory-vram-access}

The renderer reads video memory in the visible lines, and in the horizontal blanking interval too: the sprites of the
next line are fetched there (Chapter~\ref{ch:sprites}). So video memory and the sprite table are available to the CPU
only in vertical blanking, the lines after the last visible line (Chapter~\ref{ch:timing}), or while the picture is
turned off (\texttt{GD} in the Base Configuration register is 1), when the device makes no reads at all:

\begin{itemize}
	\item A write to video memory or to the sprite table at any other time is ignored, and a read returns zero. At reset
	      \texttt{GD} is 1, so a program can load video memory straight away and then clear \texttt{GD}.
	\item The palette and the registers can be read and written at any time. A palette write changes the colour of the pixels
	      drawn after it, which allows a palette change in the middle of the picture. A register write takes effect when the
	      device next samples that register, normally at the start of the next line (Chapters~\ref{ch:mode2} and~\ref{ch:mode3}).
\end{itemize}

\subsection{The Palette Boundary}
\label{sec:palette-edge}

The palette occupies 256 words, \texttt{0x6000}--\texttt{0x60FF}. The offset \texttt{0x6100} and the rest of the range
up to \texttt{0x7000} are open bus.

\section{Planning a VRAM Layout}

Nothing fixes where tile data and tilemaps live, so a program chooses. One arrangement puts the tile set at the bottom
of VRAM and the tilemaps at the top, leaving room for the tile set to grow:

\begin{table}[H]
	\centering
	\begin{tabular}{lll}
		\toprule
		\textbf{Contents}    & \textbf{Offsets}       & \textbf{Size}     \\
		\midrule
		Tile set (Mode 1)    & \texttt{0x8000} and up & 16 words per tile \\
		Background 1 tilemap & \texttt{0xF000}        & 1024 words        \\
		Background 2 tilemap & \texttt{0xF400}        & 1024 words        \\
		Background 3 tilemap & \texttt{0xF800}        & 1024 words        \\
		\bottomrule
	\end{tabular}
	\caption{An Example Mode 1 VRAM Layout}
	\label{tab:vram-layout}
\end{table}

A tilemap is a 32 by 32 table of one-word entries, so it takes 1024 words. A Mode 1 tile of 8 by 8 pixels at 4 bits per
pixel takes 16 words; a Mode 2 tile of 16 by 16 texels takes 64 words (Section~\ref{sec:m2-tile-format}). VRAM can
therefore hold at most 2048 Mode 1 tiles or 512 Mode 2 tiles, before any tilemap is placed in it.
